\section{段选择子格式}

\begin{figure}[H]
  \centering
  % figures/ch03/fig02-16.tex — auto-extracted from chapter source
\begin{tikzpicture}
  % 16 bit selector
  \foreach \i/\label in {0/RPL$_0$, 1/RPL$_1$, 2/TI, 3/Idx$_0$, 4/Idx$_1$,
                          5/Idx$_2$, 6/Idx$_3$, 7/Idx$_4$, 8/Idx$_5$,
                          9/Idx$_6$, 10/Idx$_7$, 11/Idx$_8$, 12/Idx$_9$,
                          13/Idx$_{10}$, 14/Idx$_{11}$, 15/Idx$_{12}$} {
    \node[bitcell, fill=white] at (\i*0.7, 0) {\label};
  }
  % bit numbers
  \foreach \i in {0,...,15} {
    \node[font=\tiny\ttfamily, above] at (\i*0.7+0.35, 0.3) {\i};
  }
  % grouping
  \draw[decorate, decoration={brace, amplitude=5pt, mirror}]
    (0, -0.35) -- (1.4, -0.35) node[midway, below=5pt, font=\small] {RPL (2 bit)};
  \draw[decorate, decoration={brace, amplitude=5pt, mirror}]
    (1.4, -0.35) -- (2.1, -0.35) node[midway, below=5pt, font=\small] {TI};
  \draw[decorate, decoration={brace, amplitude=5pt, mirror}]
    (2.1, -0.35) -- (11.2, -0.35) node[midway, below=5pt, font=\small] {Index (13 bit)};
\end{tikzpicture}

  \caption{段选择子的 16 位布局}
\end{figure}

\begin{bitfield}
  \bit{0--1}{RPL (Requested Privilege Level)：请求的特权级。0=内核，3=用户。
    当 selector 加载到段寄存器时，CPU 检查 $\max(\text{CPL}, \text{RPL}) \le \text{DPL}$。}
  \bit{2}{TI (Table Indicator)：0=在 GDT 中查找，1=在 LDT 中查找。我们只用 GDT。}
  \bit{3--15}{Index：GDT 中的描述符索引。GDT 条目编号从 0 开始，每个 8 字节。
    selector 的 Index 字段 $\times$ 8 就是 GDT 中的字节偏移。}
\end{bitfield}

我们的选择子取值：

\begin{longtable}{rrlll}
  \toprule
  \textbf{Selector} & \textbf{二进制} & \textbf{Index} & \textbf{RPL} & \textbf{用途} \\
  \midrule
  \hex{00} & \texttt{0000000000000000} & 0 & 0 & Null（无效，CPU 用于检查） \\
  \hex{08} & \texttt{0000000000001000} & 1 & 0 & 内核代码段 \\
  \hex{10} & \texttt{0000000000010000} & 2 & 0 & 内核数据段 \\
  \bottomrule
\end{longtable}

\codefile{src/include/arch/gdt.h}
\begin{lstlisting}[style=cstyle]
// gdt.h
enum {
    GDT_KERNEL_CODE_SEL = 0x08,  // Index=1, TI=0, RPL=0
    GDT_KERNEL_DATA_SEL = 0x10,  // Index=2, TI=0, RPL=0
};
\end{lstlisting}

% ============================================================================

